\subsection{Sequences as Conventional Interfaces}
\label{Section 2.2.3}

In working with compound data, we've stressed how data abstraction permits us
to design programs without becoming enmeshed in the details of data
representations, and how abstraction preserves for us the flexibility to
experiment with alternative representations.  In this section, we introduce
another powerful design principle for working with data structures---the use of
\newterm{conventional interfaces}.

\begin{notTaste}
Before we begin, a debt from \link{Section 2.2.1} falls due.  Every procedure
in this section builds a sequence by putting an element on the front of another
sequence, and we should say plainly what that costs us.

In the original it costs nothing.  A pair holds a pointer to the rest of the
list, so \code{cons} makes one pair and is done, whatever the length of what
follows.  Building a list of \( n \) elements this way takes \( \Theta(n) \)
operations in total.  A tuple cannot point at another tuple and call it a tail:
\code{(x,) + rest} must make a new tuple and copy every element of \code{rest}
into it.  Building up \( n \) elements therefore copies \( 1 + 2 + \dots + n \)
of them, which is \( \Theta(n^2) \).

This is not a small constant factor.  Timing the two, and doubling the length
each time:

\begin{codeBlock}
      n   tuple prepend   list append
   1000          0.69ms       0.018ms
   2000          2.49ms       0.032ms
   4000         10.09ms       0.065ms
   8000         44.86ms       0.140ms
  16000        189.54ms       0.288ms
\end{codeBlock}

\noindent
Doubling \( n \) roughly quadruples the first column and merely doubles the
second, which is what \( \Theta(n^2) \) against \( \Theta(n) \) looks like when
you watch it.  By sixteen thousand elements the difference is a factor of some
six hundred.

We are going to write the section's procedures this way regardless, because the
shape of the decomposition is what the section is about and the shape is
right.\footnote{It is worth being clear about which claims survive the change
and which do not.  The \emph{structure} of the argument---that these programs
are an enumerator, a filter, a map and an accumulator in series---is
representation-independent and holds for us exactly as it does for the
original.  The order-of-growth claims are not: where the original says a
sequence operation costs \( \Theta(n) \), ours may cost \( \Theta(n^2) \), and
we shall say so where it matters rather than quietly inheriting the assertion.}
What we should not do is pretend the cost is not there, nor conclude that
Python is unsuited to the style.  It is not: the builtins this section turns
out to duplicate are written to avoid exactly this, and by the end of the
chapter we shall have the tools to build a pair with the original's
ergonomics---one that can share a tail, and so restore the original's costs
along with its shape.
\end{notTaste}

In \link{Section 1.3} we saw how program abstractions, implemented as
higher-order procedures, can capture common patterns in programs that deal with
numerical data.  Our ability to formulate analogous operations for working with
compound data depends crucially on the style in which we manipulate our data
structures.  Consider, for example, the following procedure, analogous to the
\code{count\_leaves} procedure of \link{Section 2.2.2}, which takes a tree as
argument and computes the sum of the squares of the leaves that are odd:

\begin{codeBlock}
>>> def sum_odd_squares(tree):
...     if tree == ():
...         return 0
...     elif not is_pair(tree):
...         return square(tree) if is_odd(tree) else 0
...     else:
...         return (sum_odd_squares(tree[0])
...                 + sum_odd_squares(tree[1:]))
\end{codeBlock}

\noindent
On the surface, this procedure is very different from the following one, which
constructs a list of all the even Fibonacci numbers \( {\rm Fib}(k) \), where
\( k \) is less than or equal to a given integer \( n \):

\begin{codeBlock}
>>> def even_fibs(n):
...     def next(k):
...         if k > n:
...             return ()
...         f = fib(k)
...         if is_even(f):
...             return (f,) + next(k + 1)
...         else:
...             return next(k + 1)
...     return next(0)
\end{codeBlock}

\noindent
Despite the fact that these two procedures are structurally very different, a
more abstract description of the two computations reveals a great deal of
similarity.  The first program

\begin{itemize}
\item
enumerates the leaves of a tree;

\item
filters them, selecting the odd ones;

\item
squares each of the selected ones; and

\item
accumulates the results using \code{+}, starting with 0.
\end{itemize}

\noindent
The second program

\begin{itemize}
\item
enumerates the integers from 0 to \( n \);

\item
computes the Fibonacci number for each integer;

\item
filters them, selecting the even ones; and

\item
accumulates the results using \code{cons},  starting with the
empty list.
\end{itemize}

% figure 2.7
\begin{imageFigure}{The signal-flow plans for the procedures
  \code{sum\_odd\_squares} (top) and \code{even\_fibs} (bottom) reveal the
  commonality between the two programs.}
  \includegraphics[width=111mm]{assets/figures/chapter_2/figure_2.7d.pdf}
\end{imageFigure}

\noindent
A signal-processing engineer would find it natural to conceptualize these
processes in terms of signals flowing through a cascade of stages, each of
which implements part of the program plan, as shown in \link{Figure 2.7}.  In
\code{sum\_odd\_squares}, we begin with an \newterm{enumerator}, which
generates a ``signal'' consisting of the leaves of a given tree.  This signal
is passed through a \newterm{filter}, which eliminates all but the odd
elements.  The resulting signal is in turn passed through a \newterm{map},
which is a ``transducer'' that applies the \code{square} procedure to each
element.  The output of the map is then fed to an \newterm{accumulator}, which
combines the elements using \code{+}, starting from an initial 0.  The plan for
\code{even\_fibs} is analogous.

Unfortunately, the two procedure definitions above fail to exhibit this
signal-flow structure.  For instance, if we examine the
\code{sum\_odd\_squares} procedure, we find that the enumeration is implemented
partly by the \code{null?} and \code{pair?} tests and partly by the
tree-recursive structure of the procedure.  Similarly, the accumulation is
found partly in the tests and partly in the addition used in the recursion.  In
general, there are no distinct parts of either procedure that correspond to the
elements in the signal-flow description.  Our two procedures decompose the
computations in a different way, spreading the enumeration over the program and
mingling it with the map, the filter, and the accumulation.  If we could
organize our programs to make the signal-flow structure manifest in the
procedures we write, this would increase the conceptual clarity of the
resulting code.

\subsubsection*{Sequence Operations}

The key to organizing programs so as to more clearly reflect the signal-flow
structure is to concentrate on the ``signals'' that flow from one stage in the
process to the next.  If we represent these signals as lists, then we can use
list operations to implement the processing at each of the stages.  For
instance, we can implement the mapping stages of the signal-flow diagrams using
the \code{map} procedure from \link{Section 2.2.1}:

\begin{codeBlock}
>>> tuple(map(square, (1, 2, 3, 4, 5)))
(1, 4, 9, 16, 25)
\end{codeBlock}

\noindent
Filtering a sequence to select only those elements that satisfy a given
predicate is accomplished by

\begin{codeBlock}
>>> def filter(predicate, sequence):
...     if sequence == ():
...         return ()
...     elif predicate(sequence[0]):
...         return ((sequence[0],)
...                 + filter(predicate, sequence[1:]))
...     else:
...         return filter(predicate, sequence[1:])
\end{codeBlock}

\noindent
For example,

\begin{codeBlock}
>>> filter(is_odd, (1, 2, 3, 4, 5))
(1, 3, 5)
\end{codeBlock}

\noindent
\begin{revealTheSurface}
As with \code{map}, we may write it with a \code{for} statement:

\begin{codeBlock}
>>> def filter(predicate, sequence):
...     kept = ()
...     for item in sequence:
...         if predicate(item):
...             kept = kept + (item,)
...     return kept
\end{codeBlock}

\noindent
and as with \code{map}, Python supplies it already, lazily, so that a
\code{tuple} around it is needed to see the answer:

\begin{codeBlock}
>>> del filter
>>> tuple(filter(is_odd, (1, 2, 3, 4, 5)))
(1, 3, 5)
\end{codeBlock}

\noindent
But there is a third form, and it is the one a Python programmer reaches for.
Selecting some elements of a sequence, and transforming each element of a
sequence, are written with the same piece of syntax:

\begin{codeBlock}
>>> tuple(x for x in (1, 2, 3, 4, 5) if is_odd(x))
(1, 3, 5)
>>> tuple(square(x) for x in (1, 2, 3, 4, 5))
(1, 4, 9, 16, 25)
\end{codeBlock}

\noindent
This is a \newterm{comprehension}.  Read it from the middle outwards: \code{for
x in (1, 2, 3, 4, 5)} names each element in turn, \code{if is\_odd(x)} keeps
only some of them, and the expression at the front says what to produce from
each one that survives.  Leave off the \code{if} and nothing is filtered; make
the front expression just \code{x} and nothing is transformed.

The two stages compose in one form, which is worth noticing, because it is the
first place Python fuses two of the section's boxes into a single piece of
notation:

\begin{codeBlock}
>>> tuple(square(x) for x in (1, 2, 3, 4, 5) if is_odd(x))
(1, 9, 25)
\end{codeBlock}

\noindent
That is a filter and a map in series, written as one expression and computed in
one pass.\footnote{Written with square brackets in place of the \code{tuple(
\dots )} it is a \newterm{list comprehension}, which is the name the reader
will meet elsewhere and which produces Python's mutable \code{list}.  The form
above, in bare parentheses, is a \newterm{generator expression}: like
\code{map} and \code{filter} it is lazy, and handing it to \code{tuple} is what
draws the elements through.  We write it this way to keep our sequences
tuples.}

The accumulator has no such notation, and we shall come to why in a moment.
\end{revealTheSurface}

\noindent
Accumulations can be implemented by

\begin{codeBlock}
>>> def accumulate(op, initial, sequence):
...     if sequence == ():
...         return initial
...     else:
...         return op(sequence[0],
...                   accumulate(op, initial, sequence[1:]))

>>> accumulate(add, 0, (1, 2, 3, 4, 5))
15
>>> accumulate(mul, 1, (1, 2, 3, 4, 5))
120
>>> accumulate(lambda x, rest: (x,) + rest, (), (1, 2, 3, 4, 5))
(1, 2, 3, 4, 5)
\end{codeBlock}

\noindent
\code{add} and \code{mul} are the procedures behind \code{+} and \code{*},
which we met in \link{Exercise 1.4}; they come from the \code{operator} module,
and we use them here for the same reason the original uses \code{+} and
\code{*}---an accumulator takes a procedure, and an operator is not one until
it is named.\footnote{\code{from operator import add, mul}.  The last of the
three examples has no such name to borrow: the original passes \code{cons},
and our equivalent, joining an element to the front of a sequence, is the
\code{+} of tuples with one argument wrapped, which no module supplies
ready-made.}

\begin{revealTheSurface}
Python's name for this is \code{reduce}, from the \code{functools} module:

\begin{codeBlock}
>>> from functools import reduce
>>> reduce(add, (1, 2, 3, 4, 5), 0)
15
\end{codeBlock}

\noindent
Note that the sequence and the initial value have swapped places, and note
something less obvious that \link{Exercise 2.38} will make the whole subject
of: \code{reduce} is not quite our \code{accumulate}.  It combines in the
opposite direction.

There is no comprehension for this stage, and the reason is worth stating.  A
comprehension produces one output element for each input element it keeps; that
is what the notation is for.  An accumulation produces a single value from the
whole sequence, which is a different shape of computation, and Python spells it
with a procedure call instead.  It does supply ready-made accumulators for the
common cases---\code{sum}, \code{max}, \code{min}, \code{any}, \code{all}---and
those combine with a comprehension perfectly well, which is how most of this
section's examples are actually written in Python:

\begin{codeBlock}
>>> sum(square(x) for x in (1, 2, 3, 4, 5) if is_odd(x))
35
\end{codeBlock}
\end{revealTheSurface}

All that remains to implement signal-flow diagrams is to enumerate the sequence
of elements to be processed.  For \code{even\_fibs}, we need to generate the
sequence of integers in a given range, which we can do as follows:

\begin{codeBlock}
>>> def enumerate_interval(low, high):
...     if low > high:
...         return ()
...     else:
...         return (low,) + enumerate_interval(low + 1, high)

>>> enumerate_interval(2, 7)
(2, 3, 4, 5, 6, 7)
\end{codeBlock}

\noindent
\begin{revealTheSurface}
This one Python has as a piece of the language.  \code{range(low, high)} counts
from \code{low} up to but not including \code{high}---so the upper end is one
further along than the original's:

\begin{codeBlock}
>>> tuple(range(2, 8))
(2, 3, 4, 5, 6, 7)
\end{codeBlock}

\noindent
\code{range} is lazy like \code{map}, and it is what appears after \code{for}
in nearly every comprehension in this section.
\end{revealTheSurface}

\noindent
To enumerate the leaves of a tree, we can use\footnote{This is, in fact,
precisely the \code{fringe} procedure from \link{Exercise 2.28}.  Here we've
renamed it to emphasize that it is part of a family of general
sequence-manipulation procedures.}

\begin{codeBlock}
>>> def enumerate_tree(tree):
...     if tree == ():
...         return ()
...     elif not is_pair(tree):
...         return (tree,)
...     else:
...         return (enumerate_tree(tree[0])
...                 + enumerate_tree(tree[1:]))

>>> enumerate_tree((1, (2, (3, 4)), 5))
(1, 2, 3, 4, 5)
\end{codeBlock}

\noindent
Now we can reformulate \code{sum\_odd\_squares} and \code{even\_fibs} as in the
signal-flow diagrams.  For \code{sum\_odd\_squares}, we enumerate the sequence
of leaves of the tree, filter this to keep only the odd numbers in the
sequence, square each element, and sum the results:

\begin{codeBlock}
>>> def sum_odd_squares(tree):
...     return accumulate(
...         add, 0,
...         tuple(map(square,
...                   filter(is_odd, enumerate_tree(tree)))))
\end{codeBlock}

\noindent
For \code{even\_fibs}, we enumerate the integers from 0 to \( n \), generate
the Fibonacci number for each of these integers, filter the resulting sequence
to keep only the even elements, and accumulate the results into a list:

\begin{codeBlock}
>>> def even_fibs(n):
...     return tuple(filter(is_even,
...                         map(fib, enumerate_interval(0, n))))
\end{codeBlock}

\noindent
The second of these has lost its accumulator.  The original ends with
\code{(accumulate cons nil ...)}, which rebuilds the filtered signal into a
list---necessary there, because \code{filter} has already produced a list and
the accumulation is how the pieces are put back together.  For us the filtered
signal is a lazy sequence and \code{tuple} is what draws it out, so the final
stage is the drawing-out rather than a fold.

\begin{revealTheSurface}
Both read more directly as comprehensions, and both fuse three of the four
boxes into one expression:

\begin{codeBlock}
>>> def sum_odd_squares(tree):
...     return sum(square(x)
...                for x in enumerate_tree(tree)
...                if is_odd(x))

>>> def even_fibs(n):
...     return tuple(f for f in map(fib, range(0, n + 1))
...                  if is_even(f))
\end{codeBlock}

\noindent
It is worth being careful about what has been gained.  The signal-flow
structure has not gone away; it has been written in a notation that expresses
it in one piece instead of four nested calls.  The enumerator is the \code{for}
clause, the filter is the \code{if} clause, the map is the expression at the
front, and the accumulator is the procedure wrapped round the whole.  The
argument of this section is that recognising those four stages is what lets one
see two unlike programs as the same program---and that remains true whichever
of the two notations one writes them in.
\end{revealTheSurface}

The value of expressing programs as sequence operations is that this helps us
make program designs that are modular, that is, designs that are constructed by
combining relatively independent pieces.  We can encourage modular design by
providing a library of standard components together with a conventional
interface for connecting the components in flexible ways.

Modular construction is a powerful strategy for controlling complexity in
engineering design.  In real signal-processing applications, for example,
designers regularly build systems by cascading elements selected from
standardized families of filters and transducers.  Similarly, sequence
operations provide a library of standard program elements that we can mix and
match.  For instance, we can reuse pieces from the \code{sum\_odd\_squares} and
\code{even\_fibs} procedures in a program that constructs a list of the squares
of the first \( n + 1 \) Fibonacci numbers:

\begin{codeBlock}
>>> def list_fib_squares(n):
...     return tuple(map(square, map(fib, range(0, n + 1))))

>>> list_fib_squares(10)
(0, 1, 1, 4, 9, 25, 64, 169, 441, 1156, 3025)
\end{codeBlock}

\noindent
We can rearrange the pieces and use them in computing the product of the
squares of the odd integers in a sequence:

\begin{codeBlock}
>>> def product_of_squares_of_odd_elements(sequence):
...     return reduce(mul,
...                   (square(x) for x in sequence if is_odd(x)),
...                   1)

>>> product_of_squares_of_odd_elements((1, 2, 3, 4, 5))
225
\end{codeBlock}

\noindent
We can also formulate conventional data-processing applications in terms of
sequence operations.  Suppose we have a sequence of personnel records and we
want to find the salary of the highest-paid programmer.  Assume that we have a
selector \code{salary} that returns the salary of a record, and a predicate
\code{is\_programmer} that tests if a record is for a programmer.  Then we can
write

\begin{codeBlock}
>>> def salary_of_highest_paid_programmer(records):
...     return max((salary(r) for r in records if is_programmer(r)),
...                default=0)
\end{codeBlock}

\noindent
These examples give just a hint of the vast range of operations that can be
expressed as sequence operations.\footnote{Richard \link{Waters (1979)}
developed a program that automatically analyzes traditional Fortran programs,
viewing them in terms of maps, filters, and accumulations.  He found that fully
90 percent of the code in the Fortran Scientific Subroutine Package fits neatly
into this paradigm.  One of the reasons for the success of Lisp as a
programming language is that lists provide a standard medium for expressing
ordered collections so that they can be manipulated using higher-order
operations.  The programming language APL owes much of its power and appeal to
a similar choice. In APL all data are represented as arrays, and there is a
universal and convenient set of generic operators for all sorts of array
operations.}

Sequences, implemented here as lists, serve as a conventional interface that
permits us to combine processing modules.  Additionally, when we uniformly
represent structures as sequences, we have localized the data-structure
dependencies in our programs to a small number of sequence operations.  By
changing these, we can experiment with alternative representations of
sequences, while leaving the overall design of our programs intact.  We will
exploit this capability in \link{Section 3.5}, when we generalize the
sequence-processing paradigm to admit infinite sequences.

\begin{exerciseGroup}
\phantomsection\label{Exercise 2.33}
\exerciseHeading{Exercise 2.33:} Fill in the missing expressions
to complete the following definitions of some basic sequence-manipulation
operations as accumulations.  \code{join} is the procedure that puts an element
on the front of a sequence, \code{lambda x, rest: (x,) + rest}, which we used
above where the original passes \code{cons}; the first definition is named
\code{my\_map} so as not to shadow the built-in one:

\begin{syntaxForm}
def my_map(p, sequence):
    return accumulate(lambda x, y: \meta{??}, (), sequence)

def append(seq1, seq2):
    return accumulate(join, \meta{??}, \meta{??})

def length(sequence):
    return accumulate(\meta{??}, 0, sequence)
\end{syntaxForm}

\phantomsection\label{Exercise 2.34}
\exerciseHeading{Exercise 2.34:} Evaluating a polynomial in \( x \)
at a given value of \( x \) can be formulated as an accumulation.  We evaluate
the polynomial

$$ a_n x^n + a_{n-1} x^{n-1} + \dots + a_1 x + a_0 $$

\noindent
using a well-known algorithm called \newterm{Horner's rule}, which structures
the computation as

$$ (\dots (a_n x + a_{n-1}) x + \dots + a_1) x + a_0. $$

\noindent
In other words, we start with \( a_n \), multiply by \( x \), add \( a_{n-1}
\), multiply by \( x \), and so on, until we reach \( a_0
\).\footnote{According to \link{Knuth 1981}, this rule was formulated by W. G.
Horner early in the nineteenth century, but the method was actually used by
Newton over a hundred years earlier.  Horner's rule evaluates the polynomial
using fewer additions and multiplications than does the straightforward method
of first computing \( a_n x^n \), then adding \( a_{n-1}x^{n-1} \), and so on.
In fact, it is possible to prove that any algorithm for evaluating arbitrary
polynomials must use at least as many additions and multiplications as does
Horner's rule, and thus Horner's rule is an optimal algorithm for polynomial
evaluation.  This was proved (for the number of additions) by A. M. Ostrowski
in a 1954 paper that essentially founded the modern study of optimal
algorithms.  The analogous statement for multiplications was proved by V. Y.
Pan in 1966.  The book by \link{Borodin and Munro (1975)} provides an overview
of these and other results about optimal algorithms.}

Fill in the following template to produce a procedure that evaluates a
polynomial using Horner's rule.  Assume that the coefficients of the polynomial
are arranged in a sequence, from \( a_0 \) through \( a_n \).

\begin{syntaxForm}
def horner_eval(x, coefficient_sequence):
    return accumulate(
        lambda this_coeff, higher_terms: \meta{??},
        0,
        coefficient_sequence)
\end{syntaxForm}

For example, to compute \( 1 + 3x + 5x^3 + x^5 \) at \( x = 2 \) you
would evaluate

\begin{codeBlock}
>>> horner_eval(2, (1, 3, 0, 5, 0, 1))
\end{codeBlock}

\phantomsection\label{Exercise 2.35}
\exerciseHeading{Exercise 2.35:} Redefine \code{count\_leaves} from
\link{Section 2.2.2} as an accumulation:

\begin{syntaxForm}
def count_leaves(t):
    return accumulate(\meta{??}, \meta{??}, tuple(map(\meta{??}, \meta{??})))
\end{syntaxForm}

\phantomsection\label{Exercise 2.36}
\exerciseHeading{Exercise 2.36:} The procedure \code{accumulate\_n}
is similar to \code{accumulate} except that it takes as its third argument a
sequence of sequences, which are all assumed to have the same number of
elements.  It applies the designated accumulation procedure to combine all the
first elements of the sequences, all the second elements of the sequences, and
so on, and returns a sequence of the results.  For instance, if \code{s} is a
sequence containing four sequences, \code{((1, 2, 3), (4, 5, 6), (7, 8, 9),
(10, 11, 12))}, then the value of \code{accumulate\_n(add, 0, s)} should be the
sequence \code{(22, 26, 30)}.  Fill in the missing expressions in the following
definition of \code{accumulate\_n}:

\begin{syntaxForm}
def accumulate_n(op, init, seqs):
    if seqs[0] == ():
        return ()
    else:
        return ((accumulate(op, init, \meta{??}),)
                + accumulate_n(op, init, \meta{??}))
\end{syntaxForm}

\phantomsection\label{Exercise 2.37}
\exerciseHeading{Exercise 2.37:}
Suppose we represent vectors \( \mathbf{v} = (v_i) \) as sequences of numbers,
and matrices \( \mathbf{m} = (m_{ij}) \) as sequences of vectors (the rows of
the matrix).  For example, the matrix

$$
\left(
\begin{array}{cccc}
  1 & 2 & 3 & 4 \\
  4 & 5 & 6 & 6 \\
  6 & 7 & 8 & 9
\end{array}
\right)
$$

\noindent
is represented as the sequence \code{((1, 2, 3, 4), (4, 5, 6, 6), (6, 7, 8,
9))}.  With
this representation, we can use sequence operations to concisely express the
basic matrix and vector operations.  These operations (which are described in
any book on matrix algebra) are the following:

$$
\begin{array}{rl}
	\text{\texttt{dot\_product(v, w)}} 		& \text{returns the sum } \sum_i v_i w_i; \\
	\text{\texttt{matrix\_times\_vector(m, v)}} 	& \text{returns the vector } \mathbf{t},  \\
& \text{where } t_i = \sum_j m_{ij} v_j;      \\
\text{\texttt{matrix\_times\_matrix(m, n)}} 	& \text{returns the matrix }
\mathbf{p},  \\ & \text{where } p_{ij} = \sum_k m_{ik} n_{kj};                 
\\ \text{\texttt{transpose(m)}} 		& \text{returns the matrix } \mathbf{n},  \\
& \text{where } n_{ij} = m_{ji}.
\end{array}
$$

We can define the dot product as\footnote{This definition relies on the
built-in \code{map} taking several sequences at once, as noted in
\link{Section 2.2.1}: given two, it applies the procedure to the first element
of each, then the second of each, and so on.}

\begin{codeBlock}
>>> def dot_product(v, w):
...     return accumulate(add, 0, tuple(map(mul, v, w)))
\end{codeBlock}

Fill in the missing expressions in the following procedures for computing the
other matrix operations.  (The procedure \code{accumulate\_n} is defined in
\link{Exercise 2.36}.)

\begin{syntaxForm}
def matrix_times_vector(m, v):
    return tuple(map(\meta{??}, m))

def transpose(mat):
    return accumulate_n(\meta{??}, \meta{??}, mat)

def matrix_times_matrix(m, n):
    cols = transpose(n)
    return tuple(map(\meta{??}, m))
\end{syntaxForm}

\phantomsection\label{Exercise 2.38}
\exerciseHeading{Exercise 2.38:} The \code{accumulate} procedure
is also known as \code{fold\_right}, because it combines the first element of
the sequence with the result of combining all the elements to the right.  There
is also a \code{fold\_left}, which is similar to \code{fold\_right}, except
that it combines elements working in the opposite direction:

\begin{codeBlock}
>>> def fold_left(op, initial, sequence):
...     result = initial
...     for item in sequence:
...         result = op(result, item)
...     return result
\end{codeBlock}

\noindent
The names repay a moment's thought, because they do not mean what one first
guesses.  They do not describe the order in which the elements are visited---in
both procedures the elements are taken in the order they are written.  They
describe which \emph{end the combining starts from}, or equivalently where the
initial value sits.  Writing out both for the sequence \code{(1, 2, 3)} makes
it plain:

\begin{codeBlock}
fold_right:  op(1, op(2, op(3, initial)))
fold_left:   op(op(op(initial, 1), 2), 3)
\end{codeBlock}

\noindent
In \code{fold\_right} the initial value is at the right-hand end and the
parentheses nest rightwards, so nothing can be combined until the recursion has
run all the way to the end of the sequence and started back.  In
\code{fold\_left} the initial value is at the left-hand end, and each element
is combined with the running result as it is reached.  So both are correct at
once: \code{fold\_right} does combine the first element with the result of
combining everything to its right, as the original says, and \code{fold\_left}
does work left to right in the order written.  The first statement is about
where the nesting begins; the second is about the order of traversal, and only
\code{fold\_left}'s two agree.

Python's \code{reduce} is \code{fold\_left}.  It has no \code{fold\_right},
which is why we had to write \code{accumulate} ourselves.\footnote{This is not
an oversight.  A left fold can be written as a loop that keeps one running
value, which is cheap and never runs out of stack; a right fold must reach the
end before it can combine anything, and in Python that means recursion and a
recursion limit.  Languages that evaluate lazily prefer the right fold for
precisely the reason Python cannot: it can stop early and can work on a
sequence that has no end.}

What are the values of

\begin{codeBlock}
>>> fold_right(div, 1, (1, 2, 3))
>>> fold_left(div, 1, (1, 2, 3))
>>> fold_right(pair, (), (1, 2, 3))
>>> fold_left(pair, (), (1, 2, 3))
\end{codeBlock}

\noindent
where \code{div} is the procedure behind \code{/} and \code{pair} is
\code{lambda a, b: (a, b)}.  Work them out from the two nesting patterns above
before running them; the second pair of answers is the more instructive.

Give a property that \code{op} should satisfy to guarantee that
\code{fold\_right} and \code{fold\_left} will produce the same values for any
sequence.

\phantomsection\label{Exercise 2.39}
\exerciseHeading{Exercise 2.39:} Complete the following
definitions of \code{reverse} (\link{Exercise 2.18}) in terms of
\code{fold\_right} and \code{fold\_left} from \link{Exercise 2.38}:

\begin{syntaxForm}
def reverse(sequence):
    return fold_right(lambda x, y: \meta{??}, (), sequence)

def reverse(sequence):
    return fold_left(lambda x, y: \meta{??}, (), sequence)
\end{syntaxForm}
\end{exerciseGroup}

\subsubsection*{Nested Mappings}

We can extend the sequence paradigm to include many computations that are
commonly expressed using nested loops.\footnote{This approach to nested
mappings was shown to us by David Turner, whose languages KRC and Miranda
provide elegant formalisms for dealing with these constructs.  The examples in
this section (see also \link{Exercise 2.42}) are adapted from \link{Turner
1981}.  In \link{Section 3.5.3}, we'll see how this approach generalizes to
infinite sequences.} Consider this problem: Given a positive integer \( n \),
find all ordered pairs of distinct positive integers \( i \) and \( j \), where
\( 1 \le j < i \le n \), such that \( i + j \) is prime.  For example, if \( n
\) is 6, then the pairs are the following:

\[
\renewcommand{\arraystretch}{1.4}  % row height
\setlength{\arraycolsep}{10pt}     % space between columns
\begin{array}{c|ccccccc}
i & 2 & 3 & 4 & 4 & 5 & 6 & 6 \\
j & 1 & 2 & 1 & 3 & 2 & 1 & 5 \\
\hline
i + j & 3 & 5 & 5 & 7 & 7 & 7 & 11
\end{array}
\]

A natural way to organize this computation is to generate the sequence of all
ordered pairs of positive integers less than or equal to \( n \), filter to
select those pairs whose sum is prime, and then, for each pair \( (i, j) \)
that passes through the filter, produce the triple \( (i, j, i + j) \).

Here is a way to generate the sequence of pairs: For each integer \( i \le n
\), enumerate the integers \( j < i \), and for each such \( i \) and \( j \)
generate the pair \( (i, j) \).  In terms of sequence operations, we map along
the sequence \code{range(1, n + 1)}.  For each \( i \) in this sequence, we map
along the sequence \code{range(1, i)}.  For each \( j \) in this latter
sequence, we generate the pair \code{(i, j)}.  This gives us a sequence of
pairs for each \( i \).  Combining all the sequences for all the \( i \) (by
accumulating with \code{append}) produces the required sequence of pairs, where
\code{join\_sequences} is \code{lambda a, b: a + b}, the procedure behind the
\code{+} that runs two tuples together:

\begin{codeBlock}
>>> accumulate(
...     join_sequences, (),
...     tuple(map(lambda i: tuple(map(lambda j: (i, j),
...                                   range(1, i))),
...               range(1, n + 1))))
\end{codeBlock}

\noindent
The combination of mapping and accumulating with \code{append} is so common in
this sort of program that we will isolate it as a separate procedure:

\begin{codeBlock}
>>> def flatmap(proc, seq):
...     return accumulate(join_sequences, (),
...                       tuple(map(proc, seq)))
\end{codeBlock}

\noindent
\begin{revealTheSurface}
A comprehension can carry more than one \code{for} clause, and when it does,
the clauses nest: the second runs in full for each value the first produces,
exactly as the nested \code{map}s do.  The whole of the expression above is

\begin{codeBlock}
>>> tuple((i, j) for i in range(1, n + 1) for j in range(1, i))
\end{codeBlock}

\noindent
Read left to right, the clauses are in the same order one would write them as
nested loops: \code{for i}, then within that \code{for j}.  What has vanished
is the accumulation with \code{append}.  Nested \code{map}s produce a sequence
of sequences, which must then be flattened, and \code{flatmap} exists to
package the flattening with the mapping.  A comprehension with two \code{for}
clauses produces the flat sequence directly: there is no intermediate sequence
of sequences to collapse, because the notation never builds one.

So \code{flatmap} has no Python equivalent, and needs none---not because the
idea is missing, but because it is the ordinary reading of the notation.
\end{revealTheSurface}

Now filter this sequence of pairs to find those whose sum is prime. The filter
predicate is called for each element of the sequence; its argument is a pair
and it must extract the integers from the pair.  Thus, the predicate to apply
to each element in the sequence is

\begin{codeBlock}
>>> def is_prime_sum(pair):
...     return is_prime(pair[0] + pair[1])
\end{codeBlock}

\noindent
Finally, generate the sequence of results by mapping over the filtered pairs
using the following procedure, which constructs a triple consisting of the two
elements of the pair along with their sum:

\begin{codeBlock}
>>> def make_pair_sum(pair):
...     return (pair[0], pair[1], pair[0] + pair[1])
\end{codeBlock}

\noindent
Combining all these steps yields the complete procedure:

\begin{codeBlock}
>>> def prime_sum_pairs(n):
...     return tuple(map(
...         make_pair_sum,
...         filter(is_prime_sum,
...                flatmap(lambda i: tuple(map(lambda j: (i, j),
...                                            range(1, i))),
...                        range(1, n + 1)))))
\end{codeBlock}

\noindent
or, with the four stages in one expression:

\begin{codeBlock}
>>> def prime_sum_pairs(n):
...     return tuple((i, j, i + j)
...                  for i in range(1, n + 1)
...                  for j in range(1, i)
...                  if is_prime(i + j))

>>> prime_sum_pairs(6)
((2, 1, 3), (3, 2, 5), (4, 1, 5), (4, 3, 7), (5, 2, 7), (6, 1, 7), (6, 5, 11))
\end{codeBlock}

\noindent
Nested mappings are also useful for sequences other than those that enumerate
intervals.  Suppose we wish to generate all the permutations of a set \( S; \)
that is, all the ways of ordering the items in the set.  For instance, the
permutations of \( \{1, 2, 3\} \) are \( \{1, 2, 3\} \), \( \{1, 3, 2\} \), \(
\{2, 1, 3\} \), \( \{2, 3, 1\} \), \( \{3, 1, 2\} \), and \( \{3, 2, 1\} \).
Here is a plan for generating the permutations of \( S \): For each item \( x
\) in \( S \), recursively generate the sequence of permutations of \( S - x
\),\footnote{The set \( S - x \) is the set of all elements of \( S \),
excluding \( x \).} and adjoin \( x \) to the front of each one. This yields,
for each \( x \) in \( S \), the sequence of permutations of \( S \) that begin
with \( x \).  Combining these sequences for all \( x \) gives all the
permutations of \( S \):\footnote{The \code{\#} in Python code introduces a
\newterm{comment}.  Everything from it to the end of the line is ignored by the
interpreter; the original uses a semicolon for the same purpose.  In this book
we don't use many comments; we try to make our programs self-documenting by
using descriptive names.}

\begin{codeBlock}
>>> def permutations(s):
...     if s == ():             # empty set?
...         return ((),)        # sequence containing empty set
...     else:
...         return flatmap(
...             lambda x: tuple(map(lambda p: (x,) + p,
...                                 permutations(remove(x, s)))),
...             s)

>>> permutations((1, 2, 3))
((1, 2, 3), (1, 3, 2), (2, 1, 3), (2, 3, 1), (3, 1, 2), (3, 2, 1))
\end{codeBlock}

\noindent
Notice how this strategy reduces the problem of generating permutations of \( S
\) to the problem of generating the permutations of sets with fewer elements
than \( S \).  In the terminal case, we work our way down to the empty list,
which represents a set of no elements.  For this, we generate \code{(list
nil)}, which is a sequence with one item, namely the set with no elements.  The
\code{remove} procedure used in \code{permutations} returns all the items in a
given sequence except for a given item.  This can be expressed as a simple
filter:

\begin{codeBlock}
>>> def remove(item, sequence):
...     return tuple(x for x in sequence if x != item)
\end{codeBlock}

\begin{exerciseGroup}
\phantomsection\label{Exercise 2.40}
\exerciseHeading{Exercise 2.40:} Define a procedure
\code{unique\_pairs} that, given an integer \( n \), generates the sequence of
pairs \( (i, j) \) with \( 1 \le j < i \le n \).  Use \code{unique\_pairs}
to simplify the definition of \code{prime\_sum\_pairs} given above.

\phantomsection\label{Exercise 2.41}
\exerciseHeading{Exercise 2.41:} Write a procedure to find all
ordered triples of distinct positive integers \( i \), \( j \), and \( k \)
less than or equal to a given integer \( n \) that sum to a given integer \( s
\).

\phantomsection\label{Exercise 2.42}
\exerciseHeading{Exercise 2.42:} The ``eight-queens puzzle'' asks
how to place eight queens on a chessboard so that no queen is in check from any
other (i.e., no two queens are in the same row, column, or diagonal).  One
possible solution is shown in \link{Figure 2.8}.  One way to solve the puzzle
is to work across the board, placing a queen in each column.  Once we have
placed \( k - 1 \) queens, we must place the \( k^{\mathrm{th}} \) queen in a
position where it does not check any of the queens already on the board.  We
can formulate this approach recursively: Assume that we have already generated
the sequence of all possible ways to place \( k - 1 \) queens in the first \( k
- 1 \) columns of the board.  For each of these ways, generate an extended set
of positions by placing a queen in each row of the \( k^{\mathrm{th}} \)
column.  Now filter these, keeping only the positions for which the queen in
the \( k^{\mathrm{th}} \) column is safe with respect to the other queens.
This produces the sequence of all ways to place \( k \) queens in the first \(
k \) columns.  By continuing this process, we will produce not only one
solution, but all solutions to the puzzle.

% figure 2.8
\begin{imageFigure}{A solution to the eight-queens puzzle.}
  \includegraphics[width=48mm]{assets/figures/chapter_2/figure_2.8c.pdf}
\end{imageFigure}

We implement this solution as a procedure \code{queens}, which returns a
sequence of all solutions to the problem of placing \( n \) queens on an
\( n \times n \) chessboard.  \code{queens} has an internal procedure
\code{queen\_cols} that returns the sequence of all ways to place queens in the
first \( k \) columns of the board.\footnote{The \code{:=} in the comprehension
is an \newterm{assignment expression}: it names a value and yields it at the
same time, so that \code{positions} can be tested by \code{is\_safe} and also
produced as the result.  Without it the comprehension would have to build the
position twice, once to test and once to keep.  If it reads awkwardly---and
opinions differ---write the procedure with nested \code{for} statements
instead; the exercise is about the strategy, not the notation.}

\begin{codeBlock}

>>> def queens(board_size):
...     def queen_cols(k):
...         if k == 0:
...             return (empty_board,)
...         else:
...             return tuple(
...                 positions
...                 for rest_of_queens in queen_cols(k - 1)
...                 for new_row in range(1, board_size + 1)
...                 if is_safe(k, positions := adjoin_position(
...                     new_row, k, rest_of_queens)))
...     return queen_cols(board_size)
\end{codeBlock}

In this procedure \code{rest\_of\_queens} is a way to place \( k - 1 \) queens
in the first \( k - 1 \) columns, and \code{new\_row} is a proposed row in
which to place the queen for the \( k^{\mathrm{th}} \) column.  Complete the
program by implementing the representation for sets of board positions,
including the procedure \code{adjoin\_position}, which adjoins a new row-column
position to a set of positions, and \code{empty\_board}, which represents an
empty set of positions. You must also write the procedure \code{is\_safe},
which determines for a set of positions, whether the queen in the \(
k^{\mathrm{th}} \) column is safe with respect to the others.  (Note that we
need only check whether the new queen is safe---the other queens are already
guaranteed safe with respect to each other.)

\phantomsection\label{Exercise 2.43}
\exerciseHeading{Exercise 2.43:} Louis Reasoner is having a
terrible time doing \link{Exercise 2.42}.  His \code{queens} procedure seems to
work, but it runs extremely slowly.  (Louis never does manage to wait long
enough for it to solve even the \( 6\times6 \) case.)  When Louis asks Eva Lu
Ator for help, she points out that he has interchanged the order of the two
\code{for} clauses, writing them as

\begin{codeBlock}
...                 for new_row in range(1, board_size + 1)
...                 for rest_of_queens in queen_cols(k - 1)
\end{codeBlock}

Explain why this interchange makes the program run slowly.  Estimate how long
it will take Louis's program to solve the eight-queens puzzle, assuming that
the program in \link{Exercise 2.42} solves the puzzle in time \( T \).
\end{exerciseGroup}

